Recall the definition of normalization for deterministic transformers.
We define a deterministic transformer as normalized when it satisfies the conditions \projection\ and \concentrated.

If we make a probabilistic transformer satisfy \concentrated, then, as seen in Section \ref{prob-sim-det}, it becomes deterministic.
This changes the behavior of the transformer, so we cannot claim that for every probabilistic transformer there exists another with the same characteristics that computes the same language and satisfies \concentrated.
Thus, for probabilistic transformers, we define normalization solely as the condition \projection.


\begin{definition}[Normalized probabilistic transformer]
    We say a probabilistic transformer is normalized if it satisfies the following condition.
    \begin{itemize}
        \item \projection : The matrix of the $\OUTPUT$ has to be a projection, i.e., it is a diagonal matrix with ones on its diagonal.
    \end{itemize}
\end{definition}

Observe that the algorithm presented in Appendix \ref{appendix:alg-norm} for constructing equivalent transformers satisfying \projection\ in the deterministic setting also applies to probabilistic transformers.